\section{Conclusion}

The right way to make a hard theorem accessible is not always to flatten the full theorem at once. Sometimes the better route is to separate the core idea from the heaviest machinery and rebuild the argument in smaller pieces.

That is the philosophy of this note. The high-school warmup explains why bounded additivity forces linearity. The sphere corollary shows how orthogonality turns geometry into a one-variable identity. The matrix theorem then packages the same idea into the trace formula
\[
f(E)=\mathrm{tr}(\rho E).
\]
The result is not a replacement for Gleason's original theorem. It is a cleaner doorway into the density-matrix conclusion and the type of rigidity that makes Gleason-style results powerful in the first place.

That distinction is not a weakness to be hidden. It is the paper's organizing discipline. A reader who finishes this note should know both what has been proved and exactly what remains to be added before one returns to the original projection-only theorem. For an expository contribution, that kind of honesty is part of the result.
